\section{Actual character sections and their local Fourier cutoff limits}
\label{sec:character-quotient-repair}

The algebraic phases in MA-LA-013 do have an actual local
character-sector realization.  Ramified places require a
finite Fourier transform of the unit character, and the
resulting Gauss factor and conductor shift cannot be
replaced by an unramified Euler factor.  For odd Dirichlet
characters the real Fourier transform is not an involution;
the finite Gauss factors supply exactly the missing sign.
We construct these maps explicitly for $K=\mathbb Q$ and
for $K=\mathbb Q(\zeta_3)$, the field of MA-CH-003.
Both have class number one and finite unit group.

Let $S_f$ contain every ramified place of the finite-order
idele-class character $\beta$, and every place in the
different of $K/\mathbb Q$.  Choose a global generator
$\pi_v$ of each prime $v\in S_f$, positive when $K=\mathbb Q$.
Put
\[
\begin{aligned}
 q_v&=N\mathfrak p_v,&
 d_v&=\operatorname{ord}_v(\mathfrak d_{K_v/\mathbb Q_p}),&
 f_v&=\operatorname{ord}_v(\mathfrak f_\beta),\\
 b_v&=\beta_v(\pi_v),&
 \kappa_v&=\beta_v^{-1}|_{\mathcal O_v^\times},&
 \theta_{\beta,v}(x)&=\kappa_v(x)\mathbf1_{\mathcal O_v^\times}(x).
\end{aligned}
\tag{MA-CQ.1}
\]
Thus $|b_v|=1$.  The primitive Euler coefficient is
$c_v=b_v$ if $f_v=0$, and is zero if $f_v>0$.
Our idele convention is precisely (M22.1): the class of
the inverse prime idele has $\beta(p_v)^{-1}=c_v$ at an
unramified place.  For a primitive Dirichlet character
$\chi$, this means $c_p=\chi(p)$ and
$\beta_p|_{\mathbb Z_p^\times}=\chi_p^{-1}$ at its
conductor places.  The raw compact type is $\beta^{-1}$,
as in MA-CH-001; it is not silently inverted below.

Use the coherent local additive character $\vartheta_v$ of
MA-NR-001 (positive archimedean trace and negative finite
fractional trace) and its self-dual measure $\nu_v$.  We write the Fourier kernel as
$\overline{\vartheta_v(xy)}$, fixing every Gauss phase by this
choice.  The conventions are
\[
 \nu_{\mathbb R}=dx,\qquad
 \nu_{\mathbb C}=2\,dx\,dy,\qquad
 \nu_v(\mathcal O_v)=q_v^{-d_v/2}.
\tag{MA-CQ.2}
\]
At a finite place, the annihilator of $\mathcal O_v$ is
$\pi_v^{-d_v}\mathcal O_v$.
These measure and different normalizations are Masullo's
Definition A.3.1.2, printed p.~556, and Examples
A.3.1.5--A.3.1.10, pp.~557--558; MA-NR-001 supplies
the explicitly coherent additive signs.
The following calculation derives the phases from the
chosen Fourier integral instead of importing a root-number
formula with an unspecified sign.

\paragraph{MA-CQ-001 (the actual finite Fourier factors).}
If $f_v=0$, then
\[
 \mathcal F_v\theta_{\beta,v}(y)
   =q_v^{-d_v/2}
       \left(\mathbf1_{\pi_v^{-d_v}\mathcal O_v}(y)
         -q_v^{-1}
           \mathbf1_{\pi_v^{-(d_v+1)}\mathcal O_v}(y)\right).
\tag{MA-CQ.3}
\]
If $f_v>0$, put $m_v=f_v+d_v$ and
\[
 \tau_v(\kappa_v)
   =\int_{\mathcal O_v^\times}\kappa_v(u)
       \overline{\vartheta_v(\pi_v^{-m_v}u)}\,d\nu_v(u),
 \qquad
 \epsilon_v(\kappa_v)=q_v^{m_v/2}\tau_v(\kappa_v).
\tag{MA-CQ.4}
\]
Then
\[
\begin{aligned}
 \mathcal F_v\theta_{\beta,v}(y)
   &=\tau_v(\kappa_v)
        \kappa_v^{-1}(\pi_v^{m_v}y)
            \mathbf1_{\pi_v^{-m_v}\mathcal O_v^\times}(y),\\
 |\epsilon_v(\kappa_v)|&=1,\qquad
 \epsilon_v(\kappa_v)\epsilon_v(\kappa_v^{-1})
                                      =\kappa_v(-1).
\end{aligned}
\tag{MA-CQ.5}
\]
In particular the ramified Fourier transform is one
valuation shell, not the unramified all-pass series.

\begin{proof}
The Fourier transform of $\mathbf1_{\mathcal O_v}$ is
$q_v^{-d_v/2}\mathbf1_{\pi_v^{-d_v}\mathcal O_v}$.
Subtract the transform of $\mathbf1_{\pi_v\mathcal O_v}$,
using additive scaling by $\pi_v$, to obtain (MA-CQ.3).

For $f=f_v>0$, the function
$\kappa_v\mathbf1_{\mathcal O_v^\times}$ is invariant
under additive translations by $\pi_v^f\mathcal O_v$.
Its transform is therefore supported in the annihilator
$\pi_v^{-(f+d_v)}\mathcal O_v$.
If $\operatorname{ord}_v y>-(f+d_v)$, choose a unit
$u_0$ on which $\kappa_v$ is nontrivial but for which
$u_0-1\in\pi_v^{f-1}\mathcal O_v$.
For $f=1$ this means a suitable unit of $\mathcal O_v$;
the conductor definition supplies one.
Then $\vartheta_v(xy(u_0-1))=1$ for every unit $x$.
Changing $x$ to $u_0x$ in the Fourier integral
multiplies it by the nontrivial value $\kappa_v(u_0)$,
so that integral is zero.
Only the shell of valuation $-(f+d_v)$ remains.
For $y=\pi_v^{-(f+d_v)}v_0$ with $v_0$ a unit,
the substitution $u\mapsto v_0^{-1}u$ gives the
first formula in (MA-CQ.5).

Plancherel determines its absolute constant without
an additional Gauss-sum theorem.  The input squared
norm is $\nu_v(\mathcal O_v^\times)$.
The output squared norm is
$|\tau_v|^2q_v^{f+d_v}\nu_v(\mathcal O_v^\times)$.
Consequently $|\tau_v|=q_v^{-(f+d_v)/2}$.
For the last identity apply the Fourier transform twice.
Its square is reflection.  Scaling the shell in the
first line of (MA-CQ.5) gives
\[
 \mathcal F_v^2\theta_{\beta,v}
   =q_v^{f+d_v}
        \tau_v(\kappa_v)\tau_v(\kappa_v^{-1})
                            \theta_{\beta,v}
   =\kappa_v(-1)\theta_{\beta,v}.
\]
This proves both identities in its second line.
\end{proof}

\subsection{A section before Fourier transformation, and a Hilbert completion after it}

Let $\mu_K=\mathcal O_K^\times$, of order two or six, and
$\Gamma_S=\mu_K\prod_{v\in S_f}\pi_v^{\mathbb Z}$.
On the actual polar-free LF space
\[
 \mathcal V_S=
    \mathcal S_{\rm flat}(K_\infty)
       \mathbin{\otimes_\iota}
            \bigotimes_{v\in S_f} C_c^\infty(K_v^\times)
\tag{MA-CQ.6}
\]
let $L_\gamma f(x)=f(\gamma^{-1}x)$.
Flatness at the archimedean origin means that all real
partial derivatives vanish there.  Define
\[
 P_0=\prod_{v\in S_f}M_{\mathbf1_{\mathcal O_v^\times}},
 \qquad E_\mu=\frac1{|\mu_K|}\sum_{\zeta\in\mu_K}L_\zeta,
 \qquad
 \mathscr B_S=P_0
      \sum_{\mathbf n\in\mathbb Z^{S_f}}
             L_{\prod_v\pi_v^{n_v}}E_\mu .
\tag{MA-CQ.7}
\]
After $P_0$, the sum is finite on every LF stage.
The same proof as (MA-LA.46)--(MA-LA.47) shows that
$\mathscr B_S$ is a continuous projection with kernel
exactly the $\Gamma_S$-coinvariant relation space.
For clarity, decompose a vector into its finitely many
valuation tuples; moving each tuple to the unit tuple
changes it by a single unit relation, and averaging over
$\mu_K$ changes it by finitely many more.
The image is the unit-shell space invariant under
simultaneous multiplication by $\mu_K$.
Its kernel is closed, so this is a continuous section
of the separated coinvariant quotient.

Project further to the compact type $\beta^{-1}$ by
the actual compact-group average
\[
        \Pi_\beta f(x)=
            \int_{K_\infty^1\times\prod_{v\in S_f}\mathcal O_v^\times}
                            \beta(c)f(cx)\,dc ,
\tag{MA-CQ.8}
\]
with probability Haar measure.
This is an idempotent because the group is compact
and $\beta$ is a character.  It commutes with the
valuation section.  Its unit-shell vectors are exactly
\[
                u(x_\infty)\prod_{v\in S_f}\theta_{\beta,v}(x_v).
\tag{MA-CQ.9}
\]
For $K=\mathbb Q$, $u(-x)=(-1)^e u(x)$, where
$\chi(-1)=(-1)^e$.
For $K=\mathbb Q(\zeta_3)$ the finite-order character
is trivial at the connected complex place, so $u$
is radial.  The product in (MA-CQ.9) is invariant
under $\mu_K$, because $\beta$ is trivial on
principal ideles.  Thus this is a genuine compact
sector of the actual coinvariant quotient.

Complete this section using the archimedean norm
$L^2(K_\infty,\nu_\infty)$ on $u$, in its indicated
parity or radial subspace; call the resulting space
$\mathscr H_\beta$.  This is a specified quotient
completion, not a quotient norm induced from the
full additive $S$-adelic $L^2$ space.
The latter interpretation would contradict
(MA-LA.58).
The multiplicative norm-one compact quotient has
Haar mass one, as in MA-CH-001.
Equivalently, if a lift is integrated on the
multiplicative cover with counting measure on
$\Gamma_S$, give its compact fiber mass $|\mu_K|$;
the factor $1/|\mu_K|$ in a unit-shell lift then
gives exactly the normalized norm integral below.
This records the Haar choice independently of the
fixed self-dual additive measures (MA-CQ.2).

\paragraph{MA-CQ-002 (the actual dual-sector Fourier maps).}
On the parity or radial archimedean space define the
commuting unitaries
\[
 (V_vu)(x)=q_v^{1/2}u(\pi_v x),\qquad
 D_{v,c}=(1-\overline c\,q_v^{-1/2}V_v^*)
                              (1-cq_v^{-1/2}V_v)^{-1}.
\tag{MA-CQ.10}
\]
The archimedean rotation in $\pi_v$ acts trivially on
the complex radial sector.  Define
\[
 M_{\beta,v}=
 \begin{cases}
   b_v^{d_v}V_v^{-d_v}D_{v,\overline b_v},& f_v=0,\\
   \epsilon_v(\kappa_v)b_v^{f_v+d_v}
                                V_v^{-(f_v+d_v)},& f_v>0,
 \end{cases}
 \qquad M_{\beta,S}=\prod_{v\in S_f}M_{\beta,v}.
\tag{MA-CQ.11}
\]
Then actual Fourier transformation followed by
unit-shell periodization of (MA-CQ.9) is
\[
          H_\beta=M_{\beta,S}\mathcal F_\infty:
                        \mathscr H_\beta\longrightarrow
                                      \mathscr H_{\overline\beta}.
\tag{MA-CQ.12}
\]
It is unitary, and
\[
                         H_{\overline\beta}H_\beta=1.
\tag{MA-CQ.13}
\]
These formulas include odd parity; no assertion
$\mathcal F_{\rm odd}^2=1$ is made.
The folded formulas hold normally on compact
archimedean annuli and in the stated Hilbert norm.
The intermediate vectors need not be flat at zero.

\begin{proof}
All factors in (MA-CQ.10)--(MA-CQ.11) are unitary:
$V_v$ is normalized by the archimedean determinant
$|\pi_v|_\infty=q_v$, the numerator of $D_{v,c}$
is the adjoint of its commuting invertible
denominator, and (MA-CQ.5) gives the modulus of
each ramified scalar.

We verify that they are the actual periodization,
including the character orientation.
For $w\ne v$, $\pi_v$ is a unit at $w$.
The global character relation and the triviality
of a finite-order character on the positive
archimedean norm subgroup give
\[
                \prod_{\substack{w\in S_f\\w\ne v}}
                              \kappa_w(\pi_v)=b_v .
\tag{MA-CQ.14}
\]
After Fourier transformation the unit characters
are $\kappa_w^{-1}$.
In the term $L_{\prod_v\pi_v^{n_v}}$ of the
periodization, their angular evaluations contribute
exactly $\prod_v b_v^{n_v}$.
Indeed at $w$, removing the valuation factor
$\pi_w^{-n_w}$ leaves the unit
$\prod_{v\ne w}\pi_v^{-n_v}$, whose
$\kappa_w^{-1}$ value is
$\prod_{v\ne w}\kappa_w(\pi_v)^{n_v}$.
Multiplication over $w$ proves the assertion.

At an unramified $v$, (MA-CQ.3) permits
$n_v\le d_v$, with coefficient
$q_v^{-d_v/2}(1-q_v^{-1})$, and $n_v=d_v+1$,
with coefficient $-q_v^{-d_v/2-1}$.
Set $k=d_v-n_v\ge0$ in the first range.
The resulting archimedean operator is
\[
 b_v^{d_v}V_v^{-d_v}
 \left((1-q_v^{-1})
       \sum_{k\ge0}\overline b_v^{\,k}q_v^{-k/2}V_v^k
                   -b_v q_v^{-1/2}V_v^*\right)
       =b_v^{d_v}V_v^{-d_v}D_{v,\overline b_v}.
\tag{MA-CQ.15}
\]
The series converges absolutely in operator norm.
At a ramified $v$, the only term is
$n_v=f_v+d_v$.  Its scalar is
$\tau_v(\kappa_v)b_v^{f_v+d_v}$.
Converting the unnormalized archimedean dilation
to $V_v^{-(f_v+d_v)}$ supplies the factor
$q_v^{(f_v+d_v)/2}$, and hence gives the second
line of (MA-CQ.11).
Taking the finite product proves (MA-CQ.12).
Rapid decay of $\mathcal F_\infty u$ and its
derivatives proves normal convergence on annuli;
the geometric operator-norm bounds prove Hilbert
convergence.

This also verifies independence of a representative.
Every vector of the fixed compact type in
$\mathcal V_S$ is, modulo finitely many unit
relations, the unit-shell representative given
by (MA-CQ.7).
Additive Fourier scaling gives
$\mathcal F_S L_\gamma
       =L_{\gamma^{-1}}\mathcal F_S$
for $\gamma\in\Gamma_S$, since $|\gamma|_S=1$.
The just-established convergent folded series
can be reindexed under these translations and
kills their differences.
Thus it produces the same Hilbert vector from
every such representative.  This is a proved
extension through the chosen section, not an
application of $\mathscr B_S$ as a bounded map
on the entire additive Hilbert cover.

It remains to check Fourier square, especially
in odd parity.
On the real parity-$e$ space,
$\mathcal F_\infty^2=(-1)^e$; on the complex
radial space its square is one.
In both cases
$\mathcal F_\infty V_v\mathcal F_\infty^{-1}
        =V_v^*$.
For an unramified factor,
\[
 \mathcal F_\infty D_{v,\overline b_v}
                  \mathcal F_\infty^{-1}=D_{v,b_v}^{-1},
\]
so its factors in
$M_{\overline\beta,S}
 (\mathcal F_\infty M_{\beta,S}
                         \mathcal F_\infty^{-1})$
cancel.  At a ramified place the two scalars
multiply to $\kappa_v(-1)$ by (MA-CQ.5),
while both dilation and $b_v$ factors cancel.
Their total product is
$\prod_{v\in S_f}\kappa_v(-1)=\beta_\infty(-1)$
by the principal-idele relation.
This is $(-1)^e$ over $\mathbb Q$ and one in
the complex radial case.  Multiplication by
$\mathcal F_\infty^2$ therefore gives exactly
(MA-CQ.13).
\end{proof}

Set
\[
       D_{\beta,S}=M_{\overline\beta,S},
 \qquad Q_{\infty,Y}
    =\mathcal F_\infty^{-1}
          M_{B(|\xi|_\infty/Y)}\mathcal F_\infty .
\tag{MA-CQ.16}
\]
On the fixed archimedean parity or radial sector this also
equals $\mathcal F_\infty M_B\mathcal F_\infty^{-1}$,
because the square of $\mathcal F_\infty$ is the scalar
just computed.  The actual quotient Fourier cutoff is
\[
 H_{\overline\beta}M_{B(|\xi|_\infty/Y)}H_\beta
                  =D_{\beta,S}Q_{\infty,Y}D_{\beta,S}^* .
\tag{MA-CQ.17}
\]
This is an identity of actual bounded maps between the
specified character quotient completions.  In particular
the unramified factors on its right are $D_{v,c_v}$,
with the arithmetic coefficient $c_v$, not its
conjugate.  The intermediate Fourier map changes the
sector, which is essential to this orientation.

\subsection{The actual coupled cutoff limit, including parity and the complex place}

Let $d_\infty=1$ over $\mathbb R$ and $d_\infty=2$
over $\mathbb C$, so $|x|_\infty=|x|_{\rm Euclidean}^{d_\infty}$.
For $h\in C_c^\infty(\mathbb R_{>0})$ define
\[
 (A_hu)(x)
    =\int_0^\infty h(t)t^{-1/2}
                 u(t^{-1/d_\infty}x)\,\frac{dt}{t},
                  \qquad g(y)=h(e^y).
\tag{MA-CQ.18}
\]
This is the integrated normalized norm-dilation action
on the compact type, with norm-one Haar mass one.
The logarithmic unitary on that sector is
\[
 (\mathscr U_\infty u)(y)
   =c_\infty e^{y/2}u(e^{y/d_\infty}),
 \qquad c_{\mathbb R}=\sqrt2,\quad
                         c_{\mathbb C}=\sqrt{2\pi}.
\tag{MA-CQ.19}
\]
The complex constant uses $2\,dx\,dy$ in (MA-CQ.2).
Direct substitution gives
$\mathscr U_\infty A_h\mathscr U_\infty^{-1}=C_g$
and $\mathscr U_\infty V_v\mathscr U_\infty^{-1}
       =W_{\log q_v}$, where $(W_Rf)(y)=f(y+R)$.

Choose real profiles $A,B$ on $[0,\infty)$, smooth,
between zero and one, equal to one on $[0,1]$ and
zero on $[2,\infty)$.
For $X\ge2\epsilon>0$ and $Y>0$, put
\[
\begin{aligned}
 a_{\epsilon,X}(x)
       &=A(|x|_\infty/X)-A(|x|_\infty/\epsilon),\\
 T^{\beta,S}_{X,Y,\epsilon}(h)
       &=M_{a_{\epsilon,X}}A_h
          H_{\overline\beta}M_{B(|\xi|_\infty/Y)}H_\beta\\
       &=M_{a_{\epsilon,X}}A_h
                    D_{\beta,S}Q_{\infty,Y}D_{\beta,S}^* .
\end{aligned}
\tag{MA-CQ.20}
\]
This is an operator on the actual quotient completion
$\mathscr H_\beta$ just specified.
Its two Fourier maps act in the indicated order on
the two dual character sectors.

Write
\[
\begin{aligned}
 B_{\beta,S}(h)
   &=\sum_{\substack{v\in S_f\\ f_v=0}}\sum_{k\ge1}
      (\log q_v)q_v^{-k/2}
        \bigl(c_v^kh(q_v^k)+\overline c_v^{\,k}h(q_v^{-k})\bigr),\\
 A_e(h)&=\int_0^\infty
       \frac{e^{-(e+1/2)x}}{1-e^{-2x}}
                       \bigl(h(e^x)+h(e^{-x})\bigr)\,dx
                   \quad(K=\mathbb Q),\\
 A_{\mathbb C}(h)&=\int_0^\infty
       \frac{e^{-x/2}}{1-e^{-x}}
                       \bigl(h(e^x)+h(e^{-x})\bigr)\,dx
                   \quad(K=\mathbb Q(\zeta_3)).
\end{aligned}
\tag{MA-CQ.21}
\]
Use $A_\infty=A_e$ or $A_{\mathbb C}$ as appropriate.
For $h(1)=0$ these are ordinary convergent integrals.
The prime sums are finite.

\paragraph{MA-CQ-003 (the repaired character quotient limit).}
Every operator in (MA-CQ.20) is trace class and has a
nuclear lift into its flat unit-shell character section.
For every compact smooth $h$ with $h(1)=0$,
\[
\begin{aligned}
 \operatorname{Tr}T^{\beta,S}_{X,Y,\epsilon}(h)
       &\longrightarrow A_\infty(h)+B_{\beta,S}(h)
             &&(XY\to\infty,\ \epsilon Y\to0),\\
 \operatorname{Tr}T^{\beta,S}_{X,Y,\epsilon}(h)
       &\longrightarrow0
             &&(\epsilon Y\to\infty,\ X\ge2\epsilon).
\end{aligned}
\tag{MA-CQ.22}
\]
No A, M, or zero-location assumption is used.
For the finite-place product the actual unitary
factor has the commutator trace
\[
\begin{aligned}
 &\operatorname{Tr}\bigl(
       C_g(D_{\beta,S}^*M_\phi D_{\beta,S}-M_\phi)\bigr)\\
 &\hspace{12mm}
   =B_{\beta,S}(h)
       -\left(\sum_{v\in S_f}(f_v+d_v)\log q_v\right)g(0),
\end{aligned}
\tag{MA-CQ.23}
\]
where $\phi(-\infty)=1$, $\phi(+\infty)=0$ and
$\phi'\in C_c^\infty$.
The shift term is
$-\log(|D_K|N\mathfrak f_\beta)\,g(0)$ when $S_f$
contains all the places specified above.
It vanishes in (MA-CQ.22), but its origin in the
actual Fourier factors remains explicit.

\begin{proof}
\emph{Finite-cutoff maps and their nuclear section.}
In logarithmic coordinates $M_aA_h$ has the
smooth kernel
\[
 \bigl(A(e^y/X)-A(e^y/\epsilon)\bigr)h(e^{y-z}).
\tag{MA-CQ.24}
\]
It is supported in a compact rectangle.
Its Fourier-series rank-one expansion, with the
fixed smooth output cutoff specified in
MA-LA-011, is nuclear into the smooth compact
output interval with every derivative.
Returning by (MA-CQ.19) gives a nuclear map
into flat archimedean Schwartz functions of
the specified parity or radial type.
Multiplication by the fixed finite unit
characters places the output in the actual
section (MA-CQ.9).
The remaining factor in (MA-CQ.20) is bounded.
This proves finite-cutoff nuclearity and trace
class.  The continuous inclusion of the section
into its Hilbert completion and the same
rank-one series give the two trace pairings;
thus the quotient trace is justified by
actual maps, without extending the section
map to the full additive Hilbert cover.

\emph{The real parity and complex radial Fourier cutoffs.}
First set $D_{\beta,S}=1$ and compute the
archimedean trace directly.
Let
$K_Y=\mathcal F_\infty^{-1}(B(|\xi|_\infty/Y))$.
This is a radial Schwartz function (even in
the real case), with
\[
 K_Y(x)=Y K_1(Y^{1/d_\infty}x),
                         \qquad \int K_1\,d\nu_\infty=1.
\tag{MA-CQ.25}
\]
Extend the operator on the indicated compact
sector by zero on the other angular sectors.
The parity projector is
$\frac12(1+(-1)^e R)$; the radial projector
over $\mathbb C$ is the probability average
of rotations.
The resulting kernel is Schwartz in both
additive variables: its output variable has
compact annular support and its input
variable has uniform Schwartz decay, because
$h$ has compact norm support.
Its actual trace is therefore its diagonal
integral, justified by an absolutely
summable Fourier or Hermite rank-one
expansion as in MA-RS-002.

For the real parity-$e$ sector the trace is
\[
 \int_0^\infty h(t)t^{-1/2}
    \frac12\int_{\mathbb R}a_{\epsilon,X}(x)
       \left[
        K_Y((t^{-1}-1)x)
        +(-1)^eK_Y((t^{-1}+1)x)
       \right]dx\,\frac{dt}{t}.
\tag{MA-CQ.26}
\]
For the complex radial sector it is
\[
 \int_0^\infty h(t)t^{-1/2}
    \int_0^{2\pi}\int_{\mathbb C}
       a_{\epsilon,X}(x)
       K_Y((t^{-1/2}e^{i\theta}-1)x)
           \,d\nu_\infty(x)\,\frac{d\theta}{2\pi}\,\frac{dt}{t}.
\tag{MA-CQ.27}
\]
Both formulas are obtained from the actual
angular projector, not by assigning an
archimedean scalar distribution.

For every nonzero real or complex scalar $L$,
substitute $z=Y^{1/d_\infty}x$ in the inner
integral.  It becomes
\[
 \int\left[
        A(|z|_\infty/(XY))
         -A(|z|_\infty/(\epsilon Y))
      \right]K_1(Lz)\,d\nu_\infty(z).
\tag{MA-CQ.28}
\]
If $XY\to\infty$ and $\epsilon Y\to0$,
dominated convergence gives
$|\det_{\mathbb R}L|^{-1}$.
If $\epsilon Y\to\infty$, then also
$XY\to\infty$, and the same integral
tends to zero.
Its absolute value is bounded uniformly
by $\|K_1\|_1|\det_{\mathbb R}L|^{-1}$.
The annular multiplier is bounded by one
because $X\ge2\epsilon$ and
$a=A(|x|_\infty/X)(1-A(|x|_\infty/\epsilon))$.

For $t\ne1$ the real determinant bounds
give
\[
 \frac{t^{-1/2}}2
      \left(\frac1{|t^{-1}-1|}
                        +\frac1{t^{-1}+1}\right).
\tag{MA-CQ.29}
\]
In the complex case use the elementary
angular identity
\[
 \int_0^{2\pi}
   \frac{d\theta/(2\pi)}{|1-r e^{i\theta}|^2}
                       =\frac1{|1-r^2|}\quad(r>0,\ r\ne1).
\tag{MA-CQ.30}
\]
For $r<1$ it follows by expanding the two
geometric series and integrating their
equal powers; for $r>1$ factor out $r^2$
and apply the first case.
Thus the complex majorant after angular
integration is
$t^{-1/2}/|1-t^{-1}|$.
Multiplying either majorant by
$|h(t)|\,dt/t$ is integrable, since
$h(1)=0$ and $h$ has compact support
inside $(0,\infty)$.
The exceptional set $t=1$ has measure
zero.  Dominated convergence therefore
justifies both full trace limits.

For the first real limit the remaining
coefficient is
\[
 \frac{\sqrt t}{2}
           \left(\frac1{|1-t|}
                          +\frac{(-1)^e}{1+t}\right).
\]
Split the norm integral at $t=1$ and
substitute $t=e^x$ and $t=e^{-x}$.
The coefficient of $h(e^x)+h(e^{-x})$
is
\[
 \frac{e^{-x/2}}2
    \left(\frac1{1-e^{-x}}
                          +\frac{(-1)^e}{1+e^{-x}}\right)
       =\frac{e^{-(e+1/2)x}}{1-e^{-2x}}.
\tag{MA-CQ.31}
\]
This proves the parity formula $A_e$,
including the odd case.
The corresponding complex coefficient
is $\sqrt t/|1-t|$, which gives
$e^{-x/2}/(1-e^{-x})$ after the same
two substitutions.  This proves
$A_{\mathbb C}$ with coefficient one.
The second archimedean limit is zero.

\emph{Finite primes, conductor and different shifts.}
In logarithmic coordinates (MA-CQ.10)
has exactly the Laurent series of
MA-LA-013 with $p$ replaced by $q_v$
and phase $c_v$.  Its trace-class
piece estimate and coefficient
calculation therefore give the
$v$-term of $B_{\beta,S}$ when
$f_v=0$.
The additional factor $V_v^{-d_v}$
is a unitary shift.  Its commutator
with $M_\phi$ has multiplier
$\phi(y+d_v\log q_v)-\phi(y)$,
whose integral is $-d_v\log q_v$.
After convolution its trace is
$-d_v(\log q_v)g(0)$.
At a ramified place there is no
Euler all-pass factor: the scalar
in (MA-CQ.11) cancels from conjugation,
and its shift contributes
$-(f_v+d_v)(\log q_v)g(0)$.
All these commutators become trace
class after $C_g$, by the compact
cutoff-difference kernel estimate
of MA-LA-013.
Bounded cyclicity for the finite
commuting product gives their sum.
This proves (MA-CQ.23).
The equality
$\prod_v q_v^{f_v+d_v}
       =|D_K|N\mathfrak f_\beta$
is the norm formula for conductor
and different, with all such places
included in $S_f$.

For the exact coupled order put
$\Phi(y)=A(e^y)$ and define the fixed
trace-class operator
\[
 E_{\beta,S}(g)=
       (D_{\beta,S}^*M_\Phi D_{\beta,S}-M_\Phi)C_g.
\tag{MA-CQ.32}
\]
The order with convolution on the
right follows by taking the adjoint
of the left-order formula for
$g^*(y)=\overline{g(-y)}$.
The two directional character
coefficients in (MA-CQ.23) and
complex conjugation give the same
trace, including the real conductor
coefficient.  Hence
$\operatorname{Tr}E_{\beta,S}(g)
        =B_{\beta,S}(h)$ when $g(0)=0$.

All $D$ factors and $C_g$ commute
with $W_R$.  Positive archimedean
Fourier scaling gives on either
compact sector
\[
                W_RQ_{\infty,Y}W_{-R}
                         =Q_{\infty,e^R Y}.
\tag{MA-CQ.33}
\]
Repeating the actual bounded
cyclicity calculation (MA-LA.75),
with the right order in (MA-CQ.32),
gives
\[
\begin{aligned}
 \operatorname{Tr}T^{\beta,S}_{X,Y,\epsilon}(h)
  -\operatorname{Tr}\bigl(M_aA_hQ_{\infty,Y}\bigr)
   &=\operatorname{Tr}\bigl(E_{\beta,S}(g)Q_{\infty,XY}\bigr)\\
   &\quad-\operatorname{Tr}
           \bigl(E_{\beta,S}(g)Q_{\infty,\epsilon Y}\bigr).
\end{aligned}
\tag{MA-CQ.34}
\]
In this step $C_g$ is never commuted
past the Fourier cutoff.
The bounded multipliers
$Q_{\infty,L}$ converge strongly
to one as $L\to\infty$ and to zero
as $L\to0$.  Pairing these limits
with the fixed trace-class operator
$E_{\beta,S}(g)$ gives
$B_{\beta,S}(h)$ in the first
regime and zero in the second.
Combining with the independently
proved archimedean limits proves
(MA-CQ.22).
\end{proof}

\paragraph{MA-CQ-004 (the non-induced cubic character in an actual local quotient).}
Use the character $\psi=\psi_{269,3}$ of MA-CH-003.
Let
\[
 S_0=\{\infty,\mathfrak p_3,\mathfrak q_{269}\},
 \qquad S_1=S_0\cup\{\mathfrak p_2\},
 \qquad
 N\mathfrak p_3=3,\quad N\mathfrak q_{269}=269^2,\quad
 N\mathfrak p_2=4 .
\tag{MA-CQ.35}
\]
The section for $\beta=\psi$ consists of radial flat
complex Schwartz functions times
$\psi_{\mathfrak q_{269}}^{-1}$
on its unit shell, and the trivial unit characters
at the other selected finite places.
It is an actual section of the $S$-unit coinvariants.
Its Fourier cutoff factors can be written
\[
\begin{aligned}
 D_{\psi,S_0}
   &=\epsilon_\psi\,
       V_3^{-1}D_{3,1}\,V_{269^2}^{-1},\\
 D_{\psi,S_1}
   &=D_{\psi,S_0}D_{4,1},\qquad |\epsilon_\psi|=1 .
\end{aligned}
\tag{MA-CQ.36}
\]
Here the subscripts on $V$ name the norms of the
chosen prime generators, and
$\epsilon_\psi=\epsilon_{\mathfrak q_{269}}
       (\psi_{\mathfrak q_{269}}|_{\mathcal O^\times})$
is the explicit finite Fourier factor (MA-CQ.4).
The prime at $3$ retains its primitive Euler factor.
The prime at $269$ contributes a conductor shift
and Gauss factor, with no Euler all-pass factor.

For the test $h$ of MA-CH-004,
$h(1)=0$, $A_{\mathbb C}(h)=0$, and the actual
coupled local traces satisfy
\[
\begin{aligned}
 \lim\operatorname{Tr}T^{\psi,S_0}_{X,Y,\epsilon}(h)&=0,\\
 \lim\operatorname{Tr}T^{\psi,S_1}_{X,Y,\epsilon}(h)&=\log2
                   \qquad(XY\to\infty,\ \epsilon Y\to0).
\end{aligned}
\tag{MA-CQ.37}
\]
These are unconditional limits on the specified
local quotient completions.  They do not invoke
Hecke GRH or identify those completions with a
global narrow orthogonal projection.

\begin{proof}
The arithmetic input is the fully checked
MA-CH-003: $K=\mathbb Q(\zeta_3)$ has class
number one and unit group $\mu_6$;
$K(269^{1/3})/K$ is cyclic of degree three;
$\psi$ has conductor $(269)$.
The extension splits at the prime over $3$
and at the prime over $2$, so their primitive
coefficients are one.  The local conductor
at $\mathfrak q_{269}$ has exponent one
and its inertia invariants are zero.
The different has exponent one at
$\mathfrak p_3$ and zero at the other
two selected places.

Choose global generators $1-\zeta_3$, $269$
and $2$, respectively.  Their archimedean
norms are $3$, $269^2$ and $4$.
The uniformizer value $b_{269}$ is one:
the principal-idele relation for $269$
has trivial archimedean character and
all its other local components are
units at unramified characters.
The values $b_3=b_2=1$ are the already
verified split Frobenius values.
Substitution of these data into
(MA-CQ.11), with $\beta$ replaced by
$\overline\beta$ as in (MA-CQ.16),
gives exactly (MA-CQ.36).
The unit invariance of the section
follows from the global character
relation on $\mu_6$, and the actual
section and Fourier maps are those
proved in MA-CQ-002.
No exact numerical value of the
Gauss phase is needed; it is a
specified finite integral, with its
unit modulus proved in MA-CQ-001.

The test from MA-CH-004 has one
bump near norm four with value one
there, and a second bump in
$(1.4,1.6)$ chosen to cancel its
complex archimedean integral.
Its supports avoid the prime
powers at $3$, and there is no
Euler factor at the ramified
prime over $269$.
Thus $B_{\psi,S_0}(h)=0$.
Adding $\mathfrak p_2$ contributes
$(\log4)/\sqrt4=\log2$;
the reciprocal support contributes
zero.  The shift term in
(MA-CQ.23) is
$-\log(3\cdot269^2)h(1)=0$.
Equation (MA-CQ.22) now gives both
actual limits in (MA-CQ.37).
\end{proof}

\paragraph{Grothendieck placement and the exact global boundary.}
The finite-order character is a rank-one finite-monodromy
representation of the arithmetic fundamental group, hence
a rank-one finite-monodromy etale local system on its
unramified open arithmetic curve.
At an unramified prime its Euler polynomial is
$\det(1-\operatorname{Fr}_vT)=1-c_vT$.
At a ramified prime the determinant is taken on inertia
invariants; in the cubic example that space is zero
at $\mathfrak q_{269}$, so the primitive local factor
is one there.
These are the SGA~1, Expos\'e~V, \S7 monodromy
convention and Deligne, \emph{Valeurs de fonctions L},
Definition~6.1 and \S\S6.2--6.3, with the exact
source scope already audited in MA-CH.
For $\psi$, induction to $\mathbb Q$ is the verified
odd two-dimensional $S_3$ Artin representation of
MA-CH-003, and $\psi^c=\psi^{-1}\ne\psi$ proves that
it is not induced from a rational norm character.

The compact character projector, unit-shell section,
dual Fourier maps and Gauss factors constructed here
are analytic realizations of that character data.
The factor $q_v^{-1/2}$ and the norm-half centering
in (MA-CQ.18)--(MA-CQ.19) arise from normalized
unitary dilation and additive measure.
They are not an actual integral Tate twist of
the Artin motive.  No half-integral Tate object
or new arithmetic weight construction is being
asserted.  Finite monodromy gives the usual
weight-zero Frobenius data; it does not by itself
give a global zero-location bound for the
associated $L$-function.

All limits above concern the specifically constructed
local quotient completions.  The section is not
bounded on the full additive Hilbert cover, the
intermediate Fourier vectors need not be polar-free,
and no equality with the global weighted orthogonal
Paley--Wiener condensation is inferred.
The missing global-to-local factorization isolated
after MA-LA-013 remains a separate operator statement.
In particular neither A nor M is used here to
manufacture an independent narrow identity.
The global corrected cutoff and its discarded-divisor
criterion remain MA-WP, MA-EG and MA-PR.

\paragraph{Exact source and proof scope.}
The actual additive measures and Fourier conventions
were checked in the supplied Masullo text at Definition
A.3.1.2, p.~556, and Examples A.3.1.5--A.3.1.10,
pp.~557--558, including (A.3.1.6) and (A.3.1.8).
The different normalization is also stated at
Remark B.4.4.6 and its proof, p.~623.
The LF polar-free and finite-valuation conventions
are the explicit ones already checked in MA-LA-009--011,
with source Definition~5.3.4.1, p.~259,
pp.~260--261, and Proposition~2.1.4.1(1),
pp.~71--72.
The character convention is (M22.1)--(M22.4),
and the source compact Haar normalization is the
probability normalization of MA-CH-001.
The finite character-shell transforms, their
Gauss moduli, the parity sign, the actual
quotient Fourier maps, both archimedean
coupled cutoff limits and the finite-prime
transport are proved above.
The arithmetic of the cubic example is used
from MA-CH-003--004 with its previously
completed proof checks.
No source PDF or extraction is included in
the deliverable, and no new independent
global zero-location argument is claimed.

